\section*{Chapitre \ref{ch:iv1:calculus-and-optimisation} --- Calcul et optimisation}
\markboth{\protect\omTocMark{}{Réponses~: Calcul et optimisation}}{\protect\omTocMark{}{Réponses~: Calcul et optimisation}}

\begin{solution}{iq:iv1:calculus-and-optimisation:1}
Le minimum est atteint en $x = 1/e$, avec pour valeur $e^{-1/e} \approx 0.692$, et $e^\pi$ est le plus grand.

Écrivons $x^x = e^{x\ln x}$. La dérivée de $x\ln x$ vaut $\ln x + 1$, qui s'annule en $x = 1/e$. Lorsque $x \to 0$, la fonction tend vers 1, puis elle remonte au-delà de $1/e$~: ce point est donc le minimum.

Pour la comparaison, utilisons $g(x) = \ln x / x$\index{comparaison de puissances}. Sa dérivée $(1 - \ln x)/x^2$ est positive avant $e$ et négative après, donc $g$ atteint son maximum en $e$. D'où $\ln\pi/\pi < 1/e$, ce qui se réarrange en $e\ln\pi < \pi$, c'est-à-dire $\pi^e < e^\pi$. Les valeurs valent environ $22.46$ et $23.14$, très proches, ce qui rend une simple estimation peu sûre.

\iqwrong{Comparer les deux nombres de tête, avec des décimales approximatives. L'écart n'est que de 3~\%, si bien qu'une estimation grossière peut trancher dans un sens comme dans l'autre. Le maximum de $\ln x / x$ tranche sans aucun calcul.}

\iqfollow{Lequel est le plus grand~: $99^{100}$ ou $100^{99}$~?}{$99^{100}$. Au-delà de $e$, la fonction $\ln x / x$ décroît, donc $\ln 99/99 > \ln 100/100$, ce qui se réarrange en $100\ln 99 > 99\ln 100$.}

\iqwhy{Savoir si une fonction d'une seule ligne peut remplacer un calcul. Ramener une comparaison de puissances à la forme d'une seule fonction est le réflexe que l'on évalue.}
\end{solution}

\begin{solution}{iq:iv1:calculus-and-optimisation:2}
L'intégrale vaut $\sqrt{\pi}$, la constante de la densité normale est $1/\sqrt{2\pi}$, et $\mathbb{E}|Z| = \sqrt{2/\pi} \approx 0.798$.

Soit $I$ l'intégrale. Alors
\[
I^2 = \iint e^{-x^2 - y^2}\,dx\,dy = \int_0^{2\pi}\!\!\int_0^\infty e^{-r^2} r\,dr\,d\theta = 2\pi \cdot \tfrac12 = \pi.
\]
Le facteur $r$ issu du changement de coordonnées\index{coordonnées polaires} est ce qui rend l'intégrale intérieure exacte.

En posant $x = t/\sqrt{2}$, on obtient $\int e^{-t^2/2}\,dt = \sqrt{2\pi}$, si bien que la densité exige la constante $1/\sqrt{2\pi}$.

Enfin, par symétrie,
\[
\mathbb{E}|Z| = \frac{2}{\sqrt{2\pi}}\int_0^\infty z e^{-z^2/2}\,dz = \frac{2}{\sqrt{2\pi}} = \sqrt{\frac{2}{\pi}}.
\]
C'est inférieur à 1, comme il se doit, puisque $\mathbb{E}|Z| \le \sqrt{\mathbb{E}Z^2} = 1$.

\iqwrong{Oublier le facteur $r$, ce qui laisse une intégrale sans forme close. Ou donner $\mathbb{E}|Z| = 1$, ce qui confond la moyenne de la valeur absolue avec l'écart type.}

\iqfollow{Quelle est la valeur de $\int_0^\infty x^2 e^{-x^2}\,dx$~?}{$\sqrt{\pi}/4$. Intégrer par parties avec $u = x$, ce qui laisse la moitié de $\int_0^\infty e^{-x^2}\,dx = \sqrt{\pi}/2$.}

\iqwhy{Tout le monde cite la constante~; peu savent la retrouver. Le rapport $\sqrt{2/\pi}$ est aussi un nombre utile en pratique, reliant l'écart absolu moyen et la volatilité.}
\end{solution}

\begin{solution}{iq:iv1:calculus-and-optimisation:3}
La récurrence est $x_{n+1} = \frac12(x_n + 2/x_n)$, les itérés valent $3/2$, $17/12$ et $577/408$, et le nombre de chiffres exacts double à peu près à chaque étape.

La méthode de Newton\index{méthode de Newton} donne $x - (x^2 - 2)/(2x)$, qui se simplifie en la moyenne de $x$ et de $2/x$. C'est l'idée ancienne de moyenner une estimation trop haute avec une estimation trop basse.

À partir de 1~: $\frac12(1 + 2) = 3/2$~; puis $\frac12(3/2 + 4/3) = 17/12$~; puis $\frac12(17/12 + 24/17) = 577/408$.

Les erreurs valent environ $0.086$, $0.0025$ et $0.000\,002$. Chaque erreur vaut environ $e^2/(2x)$, où $e$ est la précédente, si bien que l'erreur est élevée au carré à chaque étape, à une constante près d'environ $0.35$.

\iqwrong{Dire que chaque étape ajoute un nombre fixe de chiffres. Cela décrit la dichotomie et les autres méthodes linéaires. Ici le nombre de chiffres exacts double.}

\iqfollow{Comment une machine sans division peut-elle calculer $1/a$~?}{Appliquer Newton à $1/x - a$, ce qui donne $x_{n+1} = x_n(2 - ax_n)$, sans division. Pour $a = 3$ à partir de $0.3$, on obtient $0.33$, puis $0.3333$. Le nombre de chiffres exacts double aussi ici.}

\iqwhy{Savoir si l'on a déjà fait l'itération à la main une fois. Le doublement des chiffres est la propriété à retenir, et le même raisonnement donne la division dans le matériel.}
\end{solution}

\begin{solution}{iq:iv1:calculus-and-optimisation:4}
$\ln 1.1 \approx 0.095$, $e^{0.1} \approx 1.105$, $\sqrt{1.04} \approx 1.020$, et l'obligation chute de 10~\% au premier ordre et de $9.5~\%$ au second ordre~; la baisse exacte est de $9.52~\%$.

Les trois développements sont $\ln(1 + x) = x - x^2/2 + x^3/3$, soit $0.1 - 0.005 + 0.000\,33 = 0.0953$~; $e^x = 1 + x + x^2/2$, soit $1.105$~; et $\sqrt{1 + x} = 1 + x/2 - x^2/8$, soit $1 + 0.02 - 0.0002 = 1.0198$.

Quel que soit le rendement initial $y$, la variation relative du prix de l'obligation est $e^{-10\,\Delta y} - 1 \approx -10\,\Delta y + \frac12(10\,\Delta y)^2$. Avec $\Delta y = 0.01$, les termes valent $-10\%$ et $+0.5\%$. Le premier est la\index{duration}duration et le second la\index{convexité}convexité. Le terme de convexité est positif quel que soit le sens des taux, ce qui explique qu'un grand mouvement laisse toujours le détenteur mieux loti que ne l'indique l'estimation par la duration.

\iqwrong{Négliger le terme de convexité et prendre $-10~\%$ pour exact. L'écart vaut environ 0.5~\% à cette taille. Pour un mouvement de deux points de pourcentage, il quadruple et atteint 2~\%.}

\iqfollow{Pourquoi la règle de 72 utilise-t-elle 72, alors que $\ln 2 = 0.693$ suggère 69~?}{Puisque $\ln(1 + r) \approx r - r^2/2$ est inférieur à $r$, le nombre exact augmente avec le taux, et 72 colle bien autour de 8~\%, où le doublement prend $9.01$ ans. Il se divise aussi facilement.}

\iqwhy{Les développements de Taylor à deux termes sont un outil quotidien du desk. Un candidat qui sait les mener vite peut aussi vérifier n'importe quel pricer.}
\end{solution}

\begin{solution}{iq:iv1:calculus-and-optimisation:5}
La première intégrale vaut $\pi/4$ pour tout $n$, et la seconde vaut $\frac{\pi}{8}\ln 2 \approx 0.272$.

Pour la première, poser $x \to \pi/2 - x$, qui échange sinus et cosinus\index{symétrie des intégrales}. En additionnant les deux écritures de l'intégrale, l'intégrande devient 1, donc $2I = \pi/2$. L'exposant ne joue aucun rôle, et c'est là tout l'intérêt.

Pour la seconde, le changement $x = \tan\theta$ élimine le dénominateur et donne
\[
I = \int_0^{\pi/4}\ln(1 + \tan\theta)\,d\theta.
\]
Poser maintenant $\theta \to \pi/4 - \theta$. La formule de soustraction de la tangente donne $1 + \tan(\pi/4 - \theta) = 2/(1 + \tan\theta)$. D'où
\[
I = \int_0^{\pi/4}\bigl(\ln 2 - \ln(1 + \tan\theta)\bigr)\,d\theta = \frac{\pi}{4}\ln 2 - I,
\]
ce qui donne $I = \frac{\pi}{8}\ln 2$.

Un contrôle d'ordre de grandeur~: l'intégrande croît de 0 à $\frac12\ln 2 \approx 0.35$ sur un intervalle de longueur 1, donc une valeur voisine de $0.27$ est plausible.

\iqwrong{Chercher une primitive. Aucune des deux intégrales n'en admet une utile. Le changement de variable qui envoie l'intervalle sur lui-même est la seule voie courte.}

\iqfollow{Calculez $\int_0^\pi \frac{x\sin x}{1 + \cos^2 x}\,dx$.}{$\pi^2/4$. La substitution $x \to \pi - x$ remplace $x$ par $\pi/2$ en moyenne, et il reste $\frac{\pi}{2}\int_0^\pi \frac{\sin x}{1 + \cos^2 x}\,dx = \frac{\pi}{2}\cdot\frac{\pi}{2}$.}

\iqwhy{La réflexion est une astuce qu'il vaut mieux avoir en tête, et le recruteur apprend à quel point le candidat s'acharne sur une primitive avant de chercher autrement.}
\end{solution}

\begin{solution}{iq:iv1:calculus-and-optimisation:6}
$n! \approx \sqrt{2\pi n}\,(n/e)^n$, donc $\binom{100}{50} \approx 1.011 \times 10^{29}$, et la probabilité d'obtenir exactement 50 piles vaut environ $0.08$.

Écrivons l'intégrande sous la forme $e^{h(x)}$ avec $h(x) = n\ln x - x$. Le maximum est atteint en $x = n$, où $h^{\prime\prime} = -n/x^2 = -1/n$. En développant à l'ordre deux\index{méthode de Laplace},
\[
n! \approx e^{n\ln n - n}\int_{-\infty}^{\infty} e^{-(x - n)^2/(2n)}\,dx = n^n e^{-n}\sqrt{2\pi n},
\]
qui est la formule de Stirling\index{formule de Stirling}. Étendre le domaine à toute la droite ne coûte presque rien, la masse se trouvant à quelques $\sqrt{n}$ de $n$. L'erreur relative vaut environ $1/(12n)$, si bien que pour $n = 10$ la formule est déjà à moins de 1~\% près.

Pour le coefficient binomial central, les puissances se simplifient~:
\[
\binom{2m}{m} \approx \frac{4^m}{\sqrt{\pi m}}.
\]
Avec $m = 50$, cela donne $1.011 \times 10^{29}$, contre la valeur exacte $1.009 \times 10^{29}$. En divisant par $2^{100}$, on obtient $1/\sqrt{50\pi} = 0.0798$, contre la valeur exacte $0.0796$.

\iqwrong{Oublier le facteur $\sqrt{2\pi n}$ et n'utiliser que $(n/e)^n$. L'erreur atteint un facteur 25 pour $n = 100$, même si cette formule donne bien le taux de croissance dominant.}

\iqfollow{Sur 1000 lancers, quelle est la probabilité d'une égalité exacte~?}{$1/\sqrt{500\pi} \approx 0.025$. Dix fois plus de lancers réduisent la probabilité d'égalité d'un facteur $\sqrt{10}$, ce que beaucoup trouvent étonnamment lent.}

\iqwhy{La démonstration est le modèle de la méthode de Laplace, qui sert à de nombreuses estimations de grandes déviations et de point-selle.}
\end{solution}

\begin{solution}{iq:iv1:calculus-and-optimisation:7}
La récurrence est $I_n = \frac{n - 1}{n}I_{n - 2}$, et $I_{10} = 63\pi/512 \approx 0.3866$. L'approximation donne $0.3963$, soit 2.5~\% de trop.

En intégrant par parties\index{intégration par parties} avec $\sin^{n - 1}x$ et $\sin x$, le terme de bord s'annule et
\[
I_n = (n - 1)\int_0^{\pi/2}\sin^{n - 2}x\cos^2 x\,dx = (n - 1)(I_{n - 2} - I_n).
\]
On résout pour obtenir la récurrence. En partant de $I_0 = \pi/2$,
\[
I_{10} = \frac{\pi}{2}\cdot\frac12\cdot\frac34\cdot\frac56\cdot\frac78\cdot\frac{9}{10} = \frac{\pi}{2}\cdot\frac{63}{256}.
\]

Pour la forme valable pour $n$ grand, la récurrence montre que $nI_nI_{n - 1}$ ne dépend pas de $n$, et qu'en $n = 1$ il vaut $\pi/2$. Comme $I_n \approx I_{n - 1}$ pour $n$ grand, on obtient $I_n \approx \sqrt{\pi/(2n)}$. L'écart en $n = 10$ devient presque nul si l'on remplace $2n$ par $2n + 1$, ce qui donne $\sqrt{\pi/21} = 0.3868$.

\iqwrong{Utiliser la récurrence avec $n/(n - 1)$, dans le mauvais sens. Les intégrales croîtraient alors, alors que $\sin^n x$ décroît avec $n$.}

\iqfollow{Quelle est la moyenne de $\sin^{10}x$ sur un cercle complet, et à quoi d'autre est-elle égale~?}{$63/256$, soit $\frac{2}{\pi}I_{10}$. C'est aussi la probabilité d'obtenir exactement cinq piles en dix lancers, $\binom{10}{5}/2^{10}$, ce qui n'est pas un hasard~: le développement de $\sin x$ en exponentielles complexes donne ce coefficient.}

\iqwhy{C'est la formule de réduction classique, et la limite de Wallis ainsi que la constante de la formule de Stirling en découlent.}
\end{solution}

\begin{solution}{iq:iv1:calculus-and-optimisation:8}
Le premier maximum vaut $1/27$ en $x = y = z = 1/3$, et le second vaut $\sqrt{2B}$ en $x = y = B/2$. Pour $B = 2$, le multiplicateur vaut $\frac12$, qui est la valeur marginale du budget.

Pour le premier, on impose que le gradient de $xyz$ soit proportionnel à celui de la contrainte\index{multiplicateurs de Lagrange}~: $yz = xz = xy = \lambda$. Les variables étant positives, on obtient $x = y = z$. L'autre voie est l'inégalité entre moyenne arithmétique et moyenne géométrique~: $(xyz)^{1/3} \le (x + y + z)/3 = 1/3$, avec égalité lorsque les trois sont égales.

Pour le second, $1/(2\sqrt{x}) = 1/(2\sqrt{y}) = \lambda$, donc $x = y = B/2$, et la valeur vaut $V(B) = \sqrt{2B}$. Pour $B = 2$~: $x = y = 1$, $V = 2$ et $\lambda = \frac12$. Par ailleurs, $V'(B) = 1/\sqrt{2B} = \frac12$. Les deux résultats concordent.

C'est le sens général du multiplicateur~: le taux auquel la meilleure valeur croît avec la contrainte, ou prix fictif\index{prix fictif}. Il indique ce que vaut la dernière unité de budget, et donc ce qu'il est raisonnable de payer pour une unité supplémentaire.

\iqwrong{Traiter le multiplicateur comme un sous-produit sans signification. Ou s'arrêter au point critique sans vérification~: ici, le coin $(B, 0)$ donne $\sqrt{2}$, inférieur à 2.}

\iqfollow{Le budget passe de 2 à $2.1$. Estimez la nouvelle meilleure valeur sans refaire la résolution.}{$2 + 0.1 \times 0.5 = 2.05$. La valeur exacte est $\sqrt{4.2} = 2.0494$. L'estimation est légèrement haute, car $V$ est concave.}

\iqwhy{Valoriser une contrainte est la fonction pratique du multiplicateur, de la limite de risque à l'allocation de capital.}
\end{solution}

\begin{solution}{iq:iv1:calculus-and-optimisation:9}
Convexes~: $-\ln x$, $x^2$, $|x|$, $\max(x, y)$. Non convexes~: $\sqrt{x}$, qui est concave, et $xy$, qui n'est ni l'un ni l'autre. Le taux composé vaut environ $8~\%$.

Les raisons, dans l'ordre. $-\ln x$ a une dérivée seconde $1/x^2 > 0$\index{convexité}. Il en va de même pour $x^2$, de dérivée seconde 2. $|x|$ vérifie l'inégalité triangulaire. $\sqrt{x}$ a une dérivée seconde négative. Un maximum de fonctions linéaires est toujours convexe. $xy$ a pour valeurs propres de la hessienne $\pm 1$, il courbe donc vers le haut le long de $x = y$ et vers le bas le long de $x = -y$.

Passons maintenant au rendement. La composition fait intervenir des logarithmes, et $\ln$ est concave, donc l'inégalité de Jensen\index{inégalité de Jensen} donne $\mathbb{E}\ln(1 + R) \le \ln(1 + \mathbb{E}R)$. En développant autour de la moyenne,
\[
\mathbb{E}\ln(1 + R) \approx \ln 1.1 - \frac{\sigma^2}{2(1.1)^2} = 0.0953 - 0.0165 = 0.079.
\]
La croissance logarithmique vaut donc environ $7.9~\%$ à cet ordre. Une intégration numérique, ainsi qu'une simulation de la loi normale, donnent $0.078$, soit $8.1~\%$ par an en taux simple. La règle usuelle $\mu - \sigma^2/2 = 8~\%$ conduit au même résultat. Sans volatilité, la formule redonne exactement 10~\%. Un rendement gaussien peut en principe descendre sous $-100~\%$, où le logarithme n'est pas défini, mais cela se situe à $5.5$ écarts types et n'affecte pas l'ordre de grandeur.

\iqwrong{Dire que $xy$ est convexe parce que chaque variable prise isolément est linéaire. Être linéaire séparément ne dit rien de la forme conjointe, comme le montre le point-selle.}

\iqfollow{Avec un rendement moyen en excès de 10~\% et une volatilité de 20~\%, quel levier maximise la croissance $L\mu - L^2\sigma^2/2$~?}{$L = \mu/\sigma^2 = 2.5$, ce qui donne une croissance de $0.125$. Avec le double de ce levier, la croissance tombe à zéro.}

\iqwhy{La traînée de volatilité est un coût direct dans toute stratégie de composition, et la concavité du logarithme en fournit à elle seule toute l'explication.}
\end{solution}

\begin{solution}{iq:iv1:calculus-and-optimisation:10}
Les pas convergents vérifient $0 < \alpha < 0.02$~; le meilleur est $\alpha = 2/101$~; il faut environ 691 pas~; Newton n'en prend qu'un~; et la quantité de mouvement de type boule lourde en demande environ 93.

Dans la base propre, les coordonnées se découplent, et chacune est multipliée par $1 - \alpha\lambda$ à chaque pas\index{descente de gradient}. Pour converger, chaque facteur doit être inférieur à 1 en valeur absolue, et la plus grande valeur propre donne $\alpha < 2/100$.

Le meilleur pas équilibre les deux facteurs~: $1 - \alpha = -(1 - 100\alpha)$, donc $\alpha = 2/101$. Les deux facteurs ont alors pour module $99/101 = 0.9802$. Pour réduire l'erreur d'un facteur $10^6$, il faut $\ln 10^6 / (-\ln 0.9802) \approx 691$ pas, ce que confirme une exécution. Un pas légèrement supérieur à $0.02$ fait croître sans limite la deuxième coordonnée.

Le nombre qui compte est le conditionnement\index{conditionnement} $\kappa = 100$~; le nombre de pas vaut environ $\frac{\kappa}{2}\ln 10^6$.

Newton multiplie le gradient par $A^{-1}$, ce qui atteint le minimum d'un coup pour toute forme quadratique. La quantité de mouvement de type boule lourde\index{quantité de mouvement}, réglée avec $\alpha = 4/121$ et $\beta = (9/11)^2$, a pour taux $(\sqrt{\kappa} - 1)/(\sqrt{\kappa} + 1) = 9/11$. Ce taux suggère 69 pas, mais la récurrence réglée a une racine double, ce qui ajoute un facteur croissant avec le nombre de pas, et une exécution en prend 93.

\iqwrong{Choisir $\alpha$ d'après la plus petite valeur propre, par exemple $\alpha = 1$. C'est excellent pour la direction plate, et cela diverge immédiatement dans la direction la plus raide.}

\iqfollow{Que fait le gradient conjugué sur cette fonction~?}{Il s'arrête en deux pas, un pour chaque valeur propre distincte, en arithmétique exacte. Une exécution directe le confirme, sans aucun réglage et avec seulement des produits matriciels.}

\iqwhy{Pourquoi les taux d'apprentissage doivent être réglés, et pourquoi c'est le mauvais conditionnement plutôt que la dimension qui fixe le coût, sont les premières leçons de l'optimisation.}
\end{solution}

\begin{solution}{iq:iv1:calculus-and-optimisation:11}
$f(x) = \operatorname{sign}(x)|x|^{1/3}$ donne $x \mapsto -2x$, ce qui s'éloigne de la racine à chaque pas. $\arctan$ échoue lorsque $|x_0| > 1.392$, et la volatilité implicite n'est sûre qu'avec un bon point de départ ou un encadrement.

Pour la racine cubique, la mise à jour vaut $x - x^{1/3}/\bigl(\frac13 x^{-2/3}\bigr) = x - 3x = -2x$. La dérivée est infinie à la racine, si bien que la tangente dépasse toujours.

Pour l'arctangente, la mise à jour vaut $x - (1 + x^2)\arctan x$. Un point de départ très éloigné a une tangente presque plate, ce qui envoie le point suivant encore plus loin. La frontière est le cycle à deux points $x \to -x$, qui résout $\arctan x = 2x/(1 + x^2)$, en $x \approx 1.392$. Les exécutions partant de $1.3$ et de $1.5$ convergent ou divergent comme prévu.

La volatilité implicite\index{volatilité implicite} a la même forme~: le prix de l'option est plat pour une faible volatilité lorsque l'option est très hors de la monnaie, ce qui donne un vega minuscule, puis le pas explose. Avec $F = 1$, $K = 1.5$, $T = 1$ et une volatilité vraie de $0.2$, un départ à $0.05$ dépasse 100 en un seul pas. Le prix est convexe en fonction de la volatilité jusqu'à un point d'inflexion, ici en $\sqrt{2\ln 1.5} = 0.90$, puis concave au-delà. Une exécution partant du point d'inflexion converge régulièrement vers la réponse.

Une méthode générale et sûre\index{Newton sécurisé} conserve un encadrement de la racine et ne prend le pas de Newton que s'il reste à l'intérieur. Sinon, elle procède par dichotomie.

\iqwrong{Supposer qu'une fonction monotone rend Newton sûr. Les deux fonctions ci-dessus sont monotones.}

\iqfollow{En dichotomie de 0.01 à 5, combien d'itérations faut-il pour atteindre une précision de $0.0001$~?}{16, puisque $\log_2(4.99/0.0001) \approx 15.6$. C'est lent mais sûr, et la borne ne dépend pas de l'option.}

\iqwhy{Les solveurs qui échouent sur des entrées rares sont à l'origine des incidents en production. Le candidat doit savoir quand Newton n'est pas sûr et disposer d'une solution de repli.}
\end{solution}

\begin{solution}{iq:iv1:calculus-and-optimisation:12}
La formule renvoie environ $7.45 \times 10^{-9}$ pour une valeur vraie de $10^{-8}$~; on utilise $c/x_{\text{grand}}$. On utilise $2\sin^2(x/2)$ pour $1 - \cos x$. On utilise une mise à jour itérative, telle que celle de Welford, pour la variance.

Dans la formule, $\sqrt{10^{16} - 4}$ est très proche de $10^8$, et l'écart entre deux flottants voisins y vaut environ $1.5 \times 10^{-8}$. La soustraction renvoie donc un petit multiple de cet écart, et la moitié d'un écart est le $7.45 \times 10^{-9}$ affiché~: une erreur de 25~\%. C'est une annulation\index{annulation en virgule flottante}~: la soustraction elle-même est exacte, mais elle révèle l'erreur déjà présente dans ses opérandes.

La correction s'appuie sur le fait que le produit des racines vaut 1. On calcule la grande racine, où les deux termes s'ajoutent, puis on divise~: $x_{\text{petit}} = 1/x_{\text{grand}}$, ce qui donne $10^{-8}$ avec toute la précision voulue.

Pour $x = 10^{-8}$, $\cos x$ s'arrondit exactement à 1, donc $1 - \cos x$ renvoie 0, alors que la valeur vraie est $5 \times 10^{-17}$. La forme en demi-angle ne comporte aucune soustraction. Des fonctions de bibliothèque comme \texttt{expm1} et \texttt{log1p} résolvent le même problème pour leurs propres cas.

Pour la variance, $\sum x^2 - (\sum x)^2/n$ soustrait deux nombres énormes et presque égaux. Avec les valeurs $10^{15} + 4$, $+7$, $+13$, $+16$, la variance vraie vaut 30, et la formule naïve est fausse de plus de 1. La mise à jour de Welford\index{algorithme de Welford} suit l'écart à la moyenne courante et renvoie 30.

\iqwrong{Accuser l'arrondi de la dernière opération. La soustraction est exacte. La faute vient de ce que les opérandes portaient déjà des erreurs, et que leurs chiffres de tête se sont annulés.}

\iqfollow{Pourquoi la somme de $10^{16}$, 1, 1, $-10^{16}$, faite dans cet ordre, donne-t-elle 0, et quelle méthode renvoie 2~?}{Chaque 1 vaut exactement la moitié de l'écart entre flottants voisins autour de $10^{16}$, et le cas limite s'arrondit vers le voisin pair, si bien que ce 1 est perdu. La sommation compensée\index{sommation compensée} de Kahan reporte la partie perdue sur le terme suivant et renvoie 2.}

\iqwhy{Les bogues numériques donnent des nombres plausibles et passent la plupart des tests. La capacité à repérer l'opération qui fait perdre de la précision est ce que la question mesure.}
\end{solution}

\begin{solution}{iq:iv1:calculus-and-optimisation:13}
Différence avant~: $h \approx 10^{-8}$, précision d'environ $10^{-8}$. Différence centrée~: $h \approx 6 \times 10^{-6}$, précision d'environ $10^{-11}$. Un pas complexe atteint une précision presque complète.

La différence avant comporte deux erreurs\index{dérivation numérique}. La troncature vaut environ $h|f^{\prime\prime}|/2$, elle diminue avec $h$. L'arrondi vaut environ $\varepsilon|f|/h$, avec $\varepsilon \approx 2.2 \times 10^{-16}$, et il augmente quand $h$ diminue. En équilibrant les deux, on obtient $h \sim \sqrt{\varepsilon} = 1.5 \times 10^{-8}$ et une erreur de même ordre de grandeur, pour une fonction et des dérivées d'ordre~1. On perd donc la moitié des chiffres, et des pas plus petits ne permettent pas de les récupérer. Un pas trop petit aggrave les choses.

La différence centrée a une troncature $h^2|f^{\prime\prime\prime}|/6$. En l'équilibrant avec l'arrondi $\varepsilon/h$, on obtient $h \sim \varepsilon^{1/3} \approx 6 \times 10^{-6}$ et une erreur d'environ $\varepsilon^{2/3} \approx 4 \times 10^{-11}$. Un balayage sur $e^x$ en $x = 1$ montre les deux minima là où on les prévoit.

Pour une fonction analytique, $f(x + ih) = f(x) + ihf'(x) + O(h^2)$, donc $\operatorname{Im} f(x + ih)/h$ donne la dérivée\index{pas complexe} sans soustraction, et donc sans annulation. Un pas de $10^{-30}$ est permis, et le résultat est juste à environ 15 chiffres près. La différentiation automatique va plus loin et n'exige pas de prolongement analytique.

\iqwrong{Dire que plus le pas est petit, mieux c'est. Avec $h = 10^{-15}$, la différence avant est surtout du bruit d'arrondi.}

\iqfollow{Quels sont le meilleur pas et la précision pour la différence seconde centrée~?}{Ici l'erreur d'arrondi vaut environ $\varepsilon/h^2$, et la troncature vaut $h^2$. On a donc $h \sim \varepsilon^{1/4} \approx 10^{-4}$, avec une précision d'environ $10^{-8}$. On perd à nouveau la moitié des chiffres.}

\iqwhy{Toute grecque calculée par perturbation présente ce compromis. Un candidat qui perturbe de $10^{-12}$ n'a jamais regardé le résultat.}
\end{solution}

\begin{solution}{iq:iv1:calculus-and-optimisation:14}
L'intégrale vaut $-\frac{\pi}{2}\ln 2 \approx -1.0888$.

Notons-la $I$. L'intégrande ne diverge en 0 que comme $\ln x$, ce qui est intégrable, donc $I$ est finie.

La réflexion $x \to \pi/2 - x$ montre que $\int_0^{\pi/2}\ln\cos x\,dx = I$\index{symétrie des intégrales}. En additionnant les deux,
\[
2I = \int_0^{\pi/2}\ln(\sin x\cos x)\,dx = \int_0^{\pi/2}\ln\frac{\sin 2x}{2}\,dx = \int_0^{\pi/2}\ln\sin 2x\,dx - \frac{\pi}{2}\ln 2.
\]
On substitue $u = 2x$ dans l'intégrale restante~:
\[
\int_0^{\pi/2}\ln\sin 2x\,dx = \frac12\int_0^{\pi}\ln\sin u\,du.
\]
Le sinus est symétrique par rapport à $\pi/2$, donc l'intégrale sur $[0, \pi]$ vaut $2I$, et le terme entier vaut $I$. Ainsi $2I = I - \frac{\pi}{2}\ln 2$, ce qui donne le résultat.

Deux vérifications. Le signe est juste, car $\ln\sin x \le 0$ sur l'intervalle. Et en divisant par la longueur $\pi/2$, on voit que la moyenne de $\ln\sin x$ vaut $-\ln 2$, donc la moyenne géométrique du sinus vaut exactement $\frac12$. Elle est inférieure à la moyenne arithmétique, $2/\pi = 0.637$, comme il se doit. L'intégration numérique donne le même nombre.

\iqwrong{S'arrêter au terme $\int\ln\sin 2x\,dx$ sans reconnaître qu'il vaut à nouveau $I$. L'auto-référence est tout le sens de l'astuce. Une seconde erreur fréquente consiste à oublier le $\frac12$ de la substitution, ce qui donne un mauvais multiple.}

\iqfollow{Calculez $\int_0^{\pi/2} x\cot x\,dx$.}{$\frac{\pi}{2}\ln 2$. On intègre par parties~: le terme de bord $x\ln\sin x$ s'annule aux deux extrémités, et il reste $-\int\ln\sin x\,dx = \frac{\pi}{2}\ln 2$.}

\iqwhy{Ramener une intégrale à elle-même, avec un coefficient différent, est une technique récurrente. Le candidat doit aussi accepter de mener une substitution qui semble ne mener nulle part.}
\end{solution}

\begin{solution}{iq:iv1:calculus-and-optimisation:15}
L'intégrale vaut $\pi/2$, avec $I(t) = \pi/2 - \arctan t$ et $I(1) = \pi/4$.

Dériver sous le signe intégral\index{dérivation sous le signe intégral} supprime le $1/x$ gênant~:
\[
I'(t) = -\int_0^\infty e^{-tx}\sin x\,dx = -\frac{1}{1 + t^2},
\]
où la dernière étape est la partie imaginaire d'une intégrale d'exponentielle complexe. Donc $I(t) = c - \arctan t$. Lorsque $t \to \infty$, $I(t) \to 0$, ce qui fixe $c = \pi/2$. Pour $t \to 0$, on obtient $\pi/2$. En $t = 1$, la valeur vaut $\pi/2 - \pi/4 = \pi/4$, et l'intégration numérique en $t = \frac12, 1, 2$ est d'accord.

La convergence vient d'une annulation. Les arcs de $\sin x/x$ alternent en signe avec une taille décroissante, donc les intégrales partielles tendent vers la limite avec une erreur d'environ $1/R$. En valeur absolue, elle diverge~: sur le $k$-ième arc, $\int|\sin x|/x\,dx$ vaut environ $2/(k\pi)$, une série harmonique qui croît comme $\frac{2}{\pi}\ln R$.

Le passage à $t = 0$ doit être justifié, car l'intégrale n'est pas absolument convergente. Le deuxième théorème de la moyenne s'applique, puisque $e^{-tx}/x$ décroît avec $x$~: le reste au-delà de $R$ est au plus $2/R$ pour tout $t \ge 0$. La convergence est donc uniforme en $t$, et $I$ est continue en 0.

\iqwrong{Appliquer directement le théorème de convergence dominée en $t = 0$. Il n'existe pas de fonction intégrable dominante, car $|\sin x|/x$ n'est pas intégrable. La borne uniforme sur le reste la remplace.}

\iqfollow{Que vaut $\int_0^\infty (\sin x/x)^2\,dx$~?}{Également $\pi/2$. Une intégration par parties la ramène à $\int_0^\infty \sin 2x / x\,dx$, qui est l'intégrale ci-dessus après un changement d'échelle. Cette fois la convergence est absolue.}

\iqwhy{L'astuce du paramètre de Feynman est fréquente dans les entretiens de quant, car elle récompense celui qui voit une voie d'accès. Le soin apporté à la limite en 0 sépare les candidats préparés des candidats rigoureux.}
\end{solution}

\begin{solution}{iq:iv1:calculus-and-optimisation:16}
L'intégrale vaut $(2\pi)^{n/2}/\sqrt{\det A}$, et la covariance vaut $A^{-1}$. Pour l'exemple, l'intégrale vaut $2\pi/\sqrt{3} \approx 3.628$, la covariance vaut $\frac13$ fois la matrice de lignes $(2, -1)$, $(-1, 2)$, et la corrélation vaut $-\frac12$.

On diagonalise $A = Q\Lambda Q^\top$ et on substitue $x = Qy$\index{intégrale gaussienne}. Comme $Q$ est orthogonale, le jacobien vaut 1, et l'exposant devient $-\frac12\sum_i\lambda_iy_i^2$. L'intégrale se factorise~:
\[
\prod_i\sqrt{\frac{2\pi}{\lambda_i}} = \frac{(2\pi)^{n/2}}{\sqrt{\det A}}.
\]

Dans les coordonnées $y$, $y_i$ est normale de variance $1/\lambda_i$, les composantes étant indépendantes. La covariance de $x$ vaut donc $Q\Lambda^{-1}Q^\top = A^{-1}$\index{matrice de précision}.

Pour l'exemple, $\det A = 3$, ce qui donne $2\pi/\sqrt{3}$~; la quadrature directe est d'accord. L'inverse est la covariance annoncée, de variance $\frac23$ et de covariance $-\frac13$, donc la corrélation vaut $-\frac12$.

Le signe est le point à expliquer. L'exposant contient le terme $-xy$, avec une entrée positive dans $A$. Ce terme pénalise le fait que $x$ et $y$ aient le même signe et récompense les signes opposés, donc une entrée positive dans la matrice de précision signifie une corrélation négative. En général, la corrélation partielle de deux variables vaut $-A_{ij}/\sqrt{A_{ii}A_{jj}}$, et une simulation de la loi normale confirme $-\frac12$.

\iqwrong{Lire la covariance directement sur $A$ et affirmer que la corrélation vaut $+\frac12$. $A$ est l'inverse de la covariance, et pour deux variables l'inversion change le signe du terme hors diagonale.}

\iqfollow{Sachant $x_2 = 1$, quels sont l'espérance et la variance de $x_1$~?}{Espérance $-\frac12$ et variance $\frac12$. La variance conditionnelle vaut $1/A_{11}$, lue directement sur la matrice de précision, et l'espérance vaut $-A_{12}/A_{11}$ fois la valeur donnée.}

\iqwhy{Les intégrales gaussiennes sont à la base de toute vraisemblance et de toute mise à jour bayésienne. La distinction entre covariance et précision est souvent le point où l'intuition d'un bon candidat se trompe.}
\end{solution}

\begin{solution}{iq:iv1:calculus-and-optimisation:17}
La solution est $x = (0.65, 0.35, 0)$ avec un niveau commun $L = 0.85$. Avec 3 unités, les trois lignes sont financées, avec $x \approx (1.37, 1.07, 0.57)$ et $L = 4.7/3 \approx 1.57$.

La dérivée de $\ln(1 + x_i/n_i) = \ln(n_i + x_i) - \ln n_i$ vaut $1/(n_i + x_i)$. Avec un multiplicateur $\nu$ associé au budget et des contraintes de non-négativité, les conditions d'optimalité sont\index{conditions de Karush--Kuhn--Tucker}
\[
\frac{1}{n_i + x_i} \le \nu, \qquad x_i \ge 0, \qquad x_i\left(\nu - \frac{1}{n_i + x_i}\right) = 0, \qquad \sum_i x_i = 1.
\]
Donc chaque ligne financée a le même $n_i + x_i = L = 1/\nu$, et $x_i = \max(L - n_i, 0)$. On parle de remplissage par niveaux\index{remplissage par niveaux} (water-filling)~: le capital est versé au-dessus des niveaux $n_i$ jusqu'à épuisement du budget.

On essaie de financer les trois~: $3L - 1.7 = 1$ donne $L = 0.9 < 1$, ce qui rend $x_3$ négatif. On finance donc deux lignes~: $2L - 0.7 = 1$ donne $L = 0.85$, puis $x_1 = 0.65$, $x_2 = 0.35$. La troisième n'est pas financée à juste titre, car sa valeur marginale en zéro, $1/1$, est inférieure à $\nu = 1/0.85$. La dernière stratégie n'est pas financée, car même sa première unité rapporte moins que le niveau courant des autres.

Avec 3 unités, $3L - 1.7 = 3$ donne $L = 4.7/3$, supérieur à 1, donc les trois lignes reçoivent désormais du capital. Un optimiseur numérique renvoie la même allocation.

\iqwrong{Résoudre le système des marginales égales sans contrainte de signe et annoncer un montant négatif, ou mettre à zéro la ligne fautive sans recalculer les autres. Le budget doit être réparti à nouveau dès qu'une ligne sort.}

\iqfollow{À partir de quel capital la troisième ligne commence-t-elle à recevoir de l'argent~?}{À $1.3$. Elle est financée dès que $L$ atteint 1, et à ce niveau les deux premières lignes détiennent $0.8 + 0.5 = 1.3$. En dessous, tout le capital va aux deux premières.}

\iqwhy{Les contraintes de signe dans les problèmes d'allocation apparaissent constamment, et les conditions KKT les rendent maniables. Le candidat qui sait vérifier la ligne non financée a compris ce que signifie le multiplicateur.}
\end{solution}

\begin{solution}{iq:iv1:calculus-and-optimisation:18}
Près d'une racine simple, l'erreur vérifie $e_{n + 1} \approx \frac{f^{\prime\prime}(r)}{2f'(r)}e_n^2$. À une racine double, elle se divise par deux à chaque pas, et l'utilisation du pas $x - 2f/f'$ rétablit une convergence rapide.

Soit $r$ la racine et $e_n = x_n - r$. La formule de Taylor en $x_n$, avec reste, donne
\[
0 = f(r) = f(x_n) - f'(x_n)e_n + \tfrac12 f^{\prime\prime}(\xi)e_n^2.
\]
On divise par $f'(x_n)$ et on réarrange. Comme $x_{n + 1} = x_n - f(x_n)/f'(x_n)$,
\[
e_{n + 1} = \frac{f^{\prime\prime}(\xi)}{2f'(x_n)}e_n^2.
\]
Si $f$ est deux fois continûment dérivable et $f'(r) \ne 0$, la fraction est bornée par un certain $M$ au voisinage de $r$. Dès que $M|e_n| < 1$, les erreurs décroissent et $|e_{n + 1}| \le M e_n^2$, ce qui définit la convergence quadratique\index{convergence quadratique}. Sur $x^2 - 2$, le rapport $e_{n + 1}/e_n^2$ tend vers $1/(2\sqrt{2})$, comme le prédit la formule.

Pour une racine double, la preuve échoue, puisque $f'(r) = 0$. Avec $f = (x - 1)^2$, le pas donne $x - (x - 1)/2$, donc $e_{n + 1} = \frac12 e_n$~: convergence linéaire de taux $\frac12$, soit un bit par pas. Une réduction de $10^6$ demande 20 pas au lieu d'environ 4. En général, pour une multiplicité $m$, le taux vaut $(m - 1)/m$.

Si $m$ est connu, on prend $x - mf/f'$. Sur l'exemple, on tombe sur la racine en un pas, et en général cela rétablit la convergence quadratique.

\iqwrong{Supposer que Newton converge toujours de façon quadratique. Ou le démontrer sans aucune hypothèse sur $f'(r)$. Le dénominateur est la condition tout entière.}

\iqfollow{Et si la multiplicité est inconnue, avec quelle précision peut-on trouver une racine double en double précision~?}{Appliquer Newton à $u = f/f'$, qui a une racine simple au même endroit. En double précision, $f$ lui-même n'est que du bruit lorsque $|x - r|$ vaut environ $10^{-8}$, si bien que seule la moitié des chiffres peut être obtenue.}

\iqwhy{Cette question distingue ceux qui se rappellent le slogan de ceux qui savent énoncer ses hypothèses. Une itération de Newton lente en pratique est souvent une racine double.}
\end{solution}

\begin{solution}{iq:iv1:calculus-and-optimisation:19}
La série géométrique donne $\sum_{k \ge 1} 1/k^2$, et l'intégrale par rotation donne $\pi^2/18 + \pi^2/9 = \pi^2/6$.

\textbf{Par la série.} En développant $1/(1 - xy) = \sum_{k \ge 0}(xy)^k$, avec des termes positifs de sorte que l'on peut permuter somme et intégrale\index{convergence monotone}, chaque terme s'intègre en $1/(k + 1)^2$.

\textbf{Par la rotation.} On pose $u = (x + y)/2$ et $v = (y - x)/2$, de sorte que $x = u - v$, $y = u + v$, le jacobien vaut 2 et $1 - xy = 1 - u^2 + v^2$. Le carré unité devient un losange avec $|v| \le u$ pour $u \le \frac12$ et $|v| \le 1 - u$ pour $u \ge \frac12$. Par symétrie en $v$, l'intégrale vaut
\[
4\int_0^{1/2}\!\!\int_0^u \frac{dv}{1 - u^2 + v^2}\,du + 4\int_{1/2}^1\!\!\int_0^{1 - u}\frac{dv}{1 - u^2 + v^2}\,du.
\]
Avec $a = \sqrt{1 - u^2}$, l'intégrale intérieure vaut $\frac1a\arctan(c/a)$, où $c$ est la borne supérieure.

Dans la première partie, $\arctan(u/\sqrt{1 - u^2}) = \arcsin u$, donc cette partie vaut
\[
4\int_0^{1/2}\frac{\arcsin u}{\sqrt{1 - u^2}}\,du = 2\left(\arcsin\tfrac12\right)^2 = \frac{\pi^2}{18}.
\]

Dans la seconde, $\arctan\sqrt{(1 - u)/(1 + u)} = \frac12\arccos u$ par la formule de l'angle moitié, donc cette partie vaut
\[
2\int_{1/2}^1\frac{\arccos u}{\sqrt{1 - u^2}}\,du = \left(\arccos\tfrac12\right)^2 = \frac{\pi^2}{9}.
\]

Les deux s'ajoutent à $\pi^2/6$, et une somme partielle directe confirme $1.6449$.

\iqwrong{Permuter somme et intégrale sans raison, ou oublier le jacobien. Le facteur 2 du changement de variables est la cause habituelle d'une constante fausse.}

\iqfollow{Quelle est la somme des $1/k^2$ pour $k$ impair, et que compte $6/\pi^2$~?}{$\pi^2/8$, car les termes pairs représentent un quart de la somme totale. Et $6/\pi^2 \approx 0.608$ est la probabilité limite que deux entiers pris au hasard soient premiers entre eux~; un dénombrement exact jusqu'à 1000 concorde à environ deux décimales.}

\iqwhy{C'est un calcul difficile mené selon un plan, et il vérifie que le candidat sait mener jusqu'au bout un changement de variables sur une région non triviale.}
\end{solution}

\begin{solution}{iq:iv1:calculus-and-optimisation:20}
La suite $d_n = \ln n! - (n + \frac12)\ln n + n$ décroît et est minorée, ce qui donne la limite, puis la formule de Wallis fixe $C = \sqrt{2\pi}$.

\textbf{Monotonie.} Par définition, $d_n - d_{n + 1} = (n + \frac12)\ln\frac{n + 1}{n} - 1$. On pose $t = 1/(2n + 1)$, de sorte que $(n + 1)/n = (1 + t)/(1 - t)$ et $n + \frac12 = 1/(2t)$. La série de $\ln\frac{1 + t}{1 - t}$ donne
\[
d_n - d_{n + 1} = \sum_{k \ge 1}\frac{t^{2k}}{2k + 1} > 0.
\]

\textbf{Borne.} En remplaçant chaque $2k + 1$ par 3, on obtient une majoration par $\frac13\cdot t^2/(1 - t^2) = \frac{1}{12n(n + 1)}$, qui vaut $\frac{1}{12n} - \frac{1}{12(n + 1)}$. Ainsi $d_n - \frac{1}{12n}$ croît avec $n$. Comme $d_1 = 1$, on obtient $d_n > \frac{11}{12}$ pour tout $n$. Une suite décroissante minorée converge, donc $n!/(\sqrt{n}(n/e)^n) \to C = e^{\lim d_n}$.

\textbf{La constante.} La formule de réduction donne
\[
I_{2n} = \frac{\pi}{2}\frac{(2n)!}{2^{2n}(n!)^2}, \qquad I_{2n + 1} = \frac{2^{2n}(n!)^2}{(2n + 1)!}.
\]
Leur rapport tend vers 1, puisque $I_n$ décroît et que $I_{2n - 1}/I_{2n + 1} = (2n + 1)/(2n) \to 1$. Ainsi\index{formule de Wallis}
\[
\frac{\pi}{2}\cdot\frac{(2n + 1)\,((2n)!)^2}{2^{4n}(n!)^4} \to 1.
\]
Substituons maintenant $n! \sim C\sqrt{n}(n/e)^n$. Toutes les puissances se simplifient, $((2n)!)^2/(2^{4n}(n!)^4) \sim 2/(C^2n)$, et la limite devient $2\pi/C^2 = 1$. Donc $C = \sqrt{2\pi}$.

\iqwrong{Utiliser le développement de Stirling plus fin pour se démontrer lui-même, ou dériver la constante par la méthode de Laplace sans contrôle de l'erreur et appeler cela une preuve. C'est une bonne heuristique, mais un recruteur qui demande la monotonie attend l'argument ci-dessus.}

\iqfollow{Quelle est la précision de la formule avec la correction $1 + 1/(12n)$ pour $n = 10$~?}{Elle donne $3\,628\,685$ contre la valeur exacte $3\,628\,800$, soit une erreur de 30 parties par million. La formule simple donne $3\,598\,696$, soit une erreur de $0.8~\%$. Le $1/(12n)$ provient de la même majoration.}

\iqwhy{Le niveau le plus difficile du chapitre. Il teste si le candidat sait rendre exacte une heuristique, et s'il sait appliquer le passage à la limite classique en gardant ses hypothèses à l'esprit.}
\end{solution}
